% Part: sets-functions-relations
% Chapter: arithmetization
% Section: cauchy

\documentclass[../../../include/open-logic-section]{subfiles}

\begin{document}

\olfileid[fa-IR]{sfr}{arith}{cauchy}

\olsection{پیوست: اعداد حقیقی به‌صورت دنباله‌های کوشی}

در \olref[cuts]{sec} اعداد حقیقی را به‌صورت برش‌های ددکیند ساختیم. در این
بخش ساختی جایگزین را توضیح می‌دهیم. این ساخت بر تعریف کوشی از چیزی که
امروزه دنبالهٔ کوشی می‌نامیم تکیه دارد؛ اما کاربرد این تعریف برای
\emph{ساختن} اعداد حقیقی از آنِ دیگر نویسندگان سدهٔ نوزدهم، به‌ویژه
وایرشتراس، هاینه، مره و کانتور است. (برای شرحی تاریخی و خواندنی، بنگرید
به \citeauthor{OConnorRobertson:RN}، \citeyear{OConnorRobertson:RN}.)

پیش از رسیدن به سدهٔ نوزدهم، شایسته است سیمون استوین (1548--1620) را در
نظر بگیریم. به‌اختصار، استوین دریافت که می‌توانیم هر عدد حقیقی را برحسب
بسط دهدهی‌اش در نظر بگیریم. بنابراین حتی عددی گنگ مانند $\sqrt{2}$ نیز
بسط دهدهیِ خوبی دارد که چنین آغاز می‌شود:
\[
	1.41421356237\ldots
\] 
مدل‌کردن بسط‌های دهدهی در نظریهٔ مجموعه‌ها بسیار آسان است: کافی است آن‌ها
را تابع‌هایی چون $d \colon \Nat \to \Nat$ در نظر بگیریم که در آن $d(n)$
رقم دهدهیِ $n$اُمی است که به آن توجه داریم. سپس به اصلاحی کوچک نیاز داریم
تا بخشی از عدد حقیقی را که پیش از ممیز دهدهی می‌آید (در اینجا فقط $1$)
در نظر بگیریم. همچنین به اصلاح دیگری (یک رابطهٔ هم‌ارزی) نیاز داریم تا
تضمین کنیم، برای نمونه، $0.999\ldots = 1$. اما عرضهٔ ساختی کاملاً دقیق از
اعداد حقیقی به شیوهٔ استوین درون نظریهٔ مجموعه‌ها دشوار نیست.

استوین موضوع اصلی ما نیست. (برای مطالب بیشتر دربارهٔ استوین، بنگرید به
\citealt{KatzKatz2012}.) اما اندیشه‌ای بسیار نزدیک به آن چنین است. به‌جای
آنکه بسط دهدهیِ $\sqrt{2}$ را مستقیماً بررسی کنیم، می‌توانیم یک
\emph{دنباله} از تقریب‌های گویای هرچه دقیق‌ترِ $\sqrt{2}$ را، از راه
بسط‌هایی هرچه دقیق‌تر، در نظر بگیریم:
\[
1, 1.4, 1.414, 1.4142, 1.41421,\ldots
\]
این اندیشه که اعداد حقیقی را می‌توان از راه «تقریب‌های هرچه بهتر» در نظر
گرفت، پایهٔ رشته‌ای دیگر از بینش‌ها را فراهم می‌کند (همانند دریافت‌هایی
که هنگام ساختن $\Rat$ از $\Int$، یا $\Int$ از $\Nat$، به کار بردیم).
بینش‌های پایه از این قرارند:
\begin{enumerate}
	\item هر عدد حقیقی را می‌توان به‌صورت بسطی دهدهی (شاید نامتناهی) نوشت.
	\item اطلاعاتِ کدگذاری‌شده در یک بسط دهدهی (شاید نامتناهی) را می‌توان
	به همان اندازه با دنباله‌ای از اعداد گویا کدگذاری کرد.
	\item دنباله‌ای از اعداد گویا را می‌توان تابعی از $\Nat$ به $\Rat$ دانست؛
	کافی است $f(n)$ را عدد گویای $n$اُم در دنباله بگیریم.
\end{enumerate}
البته نه \emph{هر} تابعی از $\Nat$ به $\Rat$ عددی حقیقی به ما می‌دهد.
برای نمونه، این تابع را در نظر بگیرید:
\[
	f(n) =\begin{cases}
			1 & \text{اگر }n\text{ فرد باشد}\\
			0 &\text{اگر }n\text{ زوج باشد}
		\end{cases}
\]
نگرانی اصلی این است که دنبالهٔ $0,1,0,1,0,1,0,\ldots$ ظاهراً به هیچ
عدد حقیقی‌ای «نزدیک نمی‌شود». پس برای آنکه دنباله‌هایی را در نظر بگیریم
که به عددی حقیقی نزدیک می‌شوند، باید توجه خود را به دنباله‌هایی محدود کنیم
که نوعی \emph{حد} دارند.

پیش‌تر در \olref[his][set][limits]{sec} با مفهوم حد روبه‌رو شده‌ایم. اما
نمی‌توانیم \emph{دقیقاً} همان تعریفی را به کار ببریم که در آنجا به کار
بردیم. عبارت «$(\forall \epsilon>0)$» در آنجا به‌طور ضمنی متضمن
کمیت‌گذاری روی اعداد حقیقی بود؛ و حدهای تابع‌هایی روی اعداد حقیقی را بررسی
می‌کردیم؛ پس استناد به آن تعریف به معنای مفروض‌گرفتن همان اعداد حقیقی‌ای
بود که دقیقاً قصد \emph{ساختن} آن‌ها را داشتیم. خوشبختانه می‌توانیم با
مفهومی بسیار نزدیک از حد کار کنیم.
\begin{defn}\ollabel{def:CauchySequence}
  تابع $f: \Nat \to \Rat$ یک \emph{دنبالهٔ کوشی} است اگر و تنها اگر برای
  هر $\epsilon \in \Rat$ مثبت داشته باشیم $(\exists \ell \in
  \Nat)(\forall m, n > \ell)|f(m) - f(n)| < \epsilon$.
\end{defn}

اندیشهٔ کلی حد همان اندیشهٔ پیشین است: اگر سطحی معین از دقت می‌خواهید
(که با~$\epsilon$ سنجیده می‌شود)، «ناحیه‌ای» برای جست‌وجو وجود دارد (هر
ورودیِ بزرگ‌تر از~$\ell$). به‌آسانی می‌توان دید که دنبالهٔ $1$، $1.4$،
$1.414$، $1.4142$، $1.41421$\ldots حد دارد: اگر بخواهید $\sqrt{2}$ را با
خطایی در حدِ $\nicefrac{1}{10^{n}}$ تقریب بزنید، کافی است به هر درایه‌ای پس
از درایهٔ $n$اُم بنگرید.

پس اندیشهٔ بدیهی این است که بگوییم هر عدد حقیقی صرفاً \emph{همان} یک
دنبالهٔ کوشی است. اما، مانند ساخت‌های $\Int$ و $\Rat$، این تصور بیش از حد
ساده‌لوحانه خواهد بود: برای هر عدد حقیقیِ معین، چندین دنبالهٔ کوشیِ
متفاوت بر آن عدد حقیقی دلالت می‌کنند. راهی ساده برای دیدن این مطلب چنین
است. اگر دنبالهٔ کوشیِ~$f$ داده شده باشد، $g$ را دقیقاً همان تابعِ~$f$
تعریف کنید، جز اینکه $g(0)\neq f(0)$. چون دو دنباله در همه‌جا پس از عدد
نخست برابرند، (در نهایت) می‌خواهیم بگوییم حد یکسانی دارند، به معنایی که
در \olref{def:CauchySequence} به کار رفت، و ازاین‌رو باید آن‌ها را
«تعریف‌کنندهٔ» همان عدد حقیقی دانست. پس واقعاً باید این دنباله‌های کوشی را
همان عدد حقیقی در نظر بگیریم.

در نتیجه، بار دیگر باید رابطه‌ای هم‌ارزی روی دنباله‌های کوشی تعریف کنیم
و اعداد حقیقی را با رابطه‌های هم‌ارزی یکی بگیریم. نخست به مفهوم تابعی نیاز
داریم که در حد به $0$ میل کند. برای هر تابع $h : \Nat \to \Rat$، می‌گوییم
\emph{$h$ به $0$ میل می‌کند} اگر و تنها اگر برای هر $\epsilon \in \Rat$
مثبت داشته باشیم $(\exists \ell \in \Nat)(\forall n > \ell)|h(n)| <
\epsilon$.\footnote{این را با تعریف $\lim_{x
\mathord{\rightarrow}\infty}f(x) = 0$ در
\olref[his][set][limits]{sec} مقایسه کنید.} افزون بر این، هرگاه $f$ و $g$
تابع‌هایی $\Nat \to \Rat$ باشند، بگذارید $(f-g)(n) = f(n) - g(n)$. اکنون
تعریف می‌کنیم:
\[
	f \Realequiv g \text{ اگر و تنها اگر $(f-g)$ به $0$ میل کند}.
\]
باید بررسی کنیم که $\Realequiv$ رابطه‌ای هم‌ارزی است؛ و چنین است. سپس اگر
بخواهیم می‌توانیم اعداد حقیقی را رده‌های هم‌ارزی، تحت $\Realequiv$، از
همهٔ دنباله‌های کوشیِ $\Nat \to \Rat$ تعریف کنیم.

\begin{prob}
بگذارید $f(n) = 0$ برای هر $n$. بگذارید $g(n) =
\frac{1}{(n+1)^2}$. نشان دهید هر دو دنبالهٔ کوشی‌اند، و در واقع حد هر دو
تابع $0$ است، پس همچنین $f \sim_\Real g$.
\end{prob}

پس از این کار، طبق معمول $\equivrep{f}{\Realequiv}$ را برای ردهٔ هم‌ارزی‌ای
می‌نویسیم که $f$ را به‌عنوان !!a{element} دارد. اما برای خوانایی، در ادامه
زیرنویس را حذف می‌کنیم و فقط $\equivrep{f}{}$ می‌نویسیم. همچنین مقرر
می‌داریم که برای هر $q \in \Rat$ داشته باشیم $q_{\Real} =
\equivrep{c_{q}}{}$، که در آن $c_{q}$ تابع ثابتِ $c_q(n) = q$ برای همهٔ
$n \in \Nat$ است. سپس روابط و اعمال پایه روی اعداد حقیقی را تعریف می‌کنیم،
برای نمونه:
\begin{align*}
	\equivrep{f}{} + \equivrep{g}{} &= 	\equivrep{(f + g)}{} \\
	\equivrep{f}{} \times 	\equivrep{g}{} &= \equivrep{(f \times g)}{} 
\end{align*}
که در آن $(f + g)(n) = f(n) + g(n)$ و $(f \times g)(n) = f(n) \times
g(n)$. البته باید بررسی کنیم که هر یک از $(f + g)$، $(f-g)$ و
$(f\times g)$، هرگاه $f$ و $g$ دنبالهٔ کوشی باشند، خود دنبالهٔ کوشی است؛
و چنین‌اند، و این کار را به شما واگذار می‌کنیم.

در پایان، مفهومی از ترتیب تعریف می‌کنیم. می‌گوییم $\equivrep{f}{}$
\emph{مثبت} است اگر و تنها اگر هم $\equivrep{f}{}\neq 0_\Rat$ و هم
$(\exists \ell \in \Nat)(\forall n > \ell)0 < f(n)$. سپس می‌گوییم
$\equivrep{f}{} < \equivrep{g}{}$ اگر و تنها اگر
$\equivrep{(g - f)}{}$ مثبت باشد. باید بررسی کنیم که این تعریف خوش‌تعریف
است (یعنی به انتخاب تابع «نماینده» از ردهٔ هم‌ارزی وابسته نیست).
%\begin{align*}
%	\equivrep{f}{} \leq \equivrep{g}{} \text{ iff }&\equivrep{f}{} = \equivrep{g}{} \text{ or }\\
%	&(\exists \ell \in \Nat)(\forall n \in \Nat)(\ell < n \lif f(n) \leq g(n))
%\end{align*}
اما پس از انجام این کار، به‌آسانی می‌توان نشان داد که این تعریف‌ها خواص
جبریِ درست را به دست می‌دهند؛ یعنی:
\begin{thm}\ollabel{thm:cauchyorderedfield}
دنباله‌های کوشی یک میدان مرتب تشکیل می‌دهند.
\end{thm}

\begin{proof}
تمرین.
\end{proof}

\begin{prob}
ثابت کنید دنباله‌های کوشی یک میدان مرتب تشکیل می‌دهند.
\end{prob}

اثبات اینکه اعداد حقیقیِ بدین‌گونه ساخته‌شده خاصیت تمامیت دارند دشوارتر
است؛ پس برهان را عرضه می‌کنیم.

\begin{thm}
هر مجموعهٔ ناتهی از دنباله‌های کوشی که کران بالا داشته باشد، کوچک‌ترین
کران بالا دارد.
\end{thm}

\begin{proof}[طرح برهان] بگذارید $S$ هر مجموعهٔ ناتهی از دنباله‌های کوشی
باشد که کران بالا دارد. پس $p \in \Rat$ای هست که $p_{\Real}$ کران بالای
$S$ است. بگذارید $r \in S$؛ آنگاه $q \in \Rat$ای هست چنان‌که
$q_{\Real} < r$. پس اگر کوچک‌ترین کران بالای $S$ وجود داشته باشد، میان
$q_\Real$ و $p_\Real$ (با احتساب دو سر) قرار دارد.

با نزدیک‌شدن هم‌زمان از پایین و بالا، به ک.ک.ب. نزدیک خواهیم شد. به‌ویژه،
دو تابع $f, g \colon \Nat \to \Rat$ را تعریف می‌کنیم، با این هدف که $f$
از بالا به ک.ک.ب.\ نزدیک شود و $g$ از پایین به آن نزدیک شود. کار را با
تعریف زیر آغاز می‌کنیم:
\begin{align*}
	f(0) &= p \\
	g(0) &= q
\end{align*}
سپس، با $a_n = \frac{f(n) + g(n)}{2}$، بگذارید:\footnote{این تعریفی
بازگشتی است. اما \emph{هنوز} هیچ دلیلی عرضه نکرده‌ایم که تعریف‌های
بازگشتی مجازند.}
\begin{align*}
	f(n+1) &=
	\begin{cases}
		a_n &\text{اگر }(\forall h \in S)\equivrep{h}{} \leq (a_n)_\Real\\
		f(n)&\text{در غیر این صورت}
	\end{cases}\\
	g(n+1) &=
	\begin{cases}
		a_n &\text{اگر }(\exists h \in S)\equivrep{h}{} \geq (a_n)_\Real\\
	 	g(n) &\text{در غیر این صورت}
	\end{cases}
\end{align*}
هر دو $f$ و $g$ دنبالهٔ کوشی‌اند. (این را نسبتاً آسان می‌توان بررسی کرد،
اما به‌عنوان تمرین واگذار می‌کنیم.) توجه کنید تابع $(f-g)$ به $0$ میل
می‌کند، زیرا اختلاف $f$ و $g$ در هر گام نصف می‌شود. پس
$\equivrep{f}{} = \equivrep{g}{}$.

نخست نشان می‌دهیم $\equivrep{f}{}$ کران بالای $S$ است، یعنی\ $(\forall h \in S)\equivrep{h}{} \leq \equivrep{f}{}$.
(در طول کار به \olref{thm:cauchyorderedfield} استناد می‌کنیم.) بگذارید
$h \in S$ و برای برهان خلف فرض کنید $\equivrep{f}{} < \equivrep{h}{}$،
پس $0_\Real < \equivrep{(h-f)}{}$. چون $f$ دنبالهٔ کوشیِ یکنواخت نزولی
است، $n \in \Nat$ای وجود دارد چنان‌که
$\equivrep{(c_{f(n)} - f)}{} < \equivrep{(h-f)}{}$. بنابراین:
\[
	(f(n))_\Real = \equivrep{c_{f(n)}}{} < \equivrep{f}{} + \equivrep{(h-f)}{} = \equivrep{h}{},
\]
که با این واقعیت تناقض دارد که بنا به ساخت، $\equivrep{h}{} \leq (f(n))_\Real$.

سپس نشان می‌دهیم $\equivrep{f}{} = \equivrep{g}{}$ \emph{کوچک‌ترین}
کران بالای $S$ است. پس بگذارید $j$ هر دنبالهٔ کوشی باشد و فرض کنید
$\equivrep{j}{} < \equivrep{g}{}$. با استدلالی مانند بالا (و استفاده از
این واقعیت که $g$ \emph{صعودی} است)، $n \in \Nat$ای وجود دارد چنان‌که
$\equivrep{j}{} < (g(n))_\Real$. اما بنا به ساخت، $h \in S$ای وجود دارد
چنان‌که $(g(n))_\Real \leq \equivrep{h}{}$؛ پس $\equivrep{j}{} <
\equivrep{h}{}$ و بنابراین $\equivrep{j}{}$ کران بالای $S$ نیست.
\end{proof}

\end{document}
